\documentclass[11pt,a4paper]{article}
\usepackage[T1]{fontenc}
\usepackage[utf8]{inputenc}
\usepackage{lmodern,amsmath,amssymb}
\usepackage[margin=25mm]{geometry}
\usepackage{longtable,booktabs,array,calc}
\usepackage[hidelinks]{hyperref}
\hypersetup{pdftitle={The mass gap: current position of this programme},pdfauthor={navstokgap research working notes}}
\newcommand{\tightlist}{\setlength{\itemsep}{0pt}\setlength{\parskip}{0pt}}
\setlength{\emergencystretch}{3em}
\setcounter{secnumdepth}{0}
\begin{document}
\hypertarget{the-mass-gap-current-position-of-this-programme}{%
\section{The mass gap: current position of this
programme}\label{the-mass-gap-current-position-of-this-programme}}

This is the synthesis of the notes written since the goal was set to a
proof of the Yang--Mills existence and mass-gap conjecture. It states
what has been proved here, what has been imported, what has been ruled
out, and what remains, with constants explicit and with each claim
labelled. The short version: the conjecture decomposes into six named
statements, of which two are proved and one is imported; the
\textbf{upper side of the problem, the clause \(m<\infty\), is closed}
as far as any variational method can take it, and reduced to the
existence of the theory plus the nontriviality of one flowed correlator;
the \textbf{lower side, the clause \(m>0\), is untouched}, and three
natural routes to it have been closed by explicit computation rather
than by difficulty. What remains is a controlled truncation at each
scale, which is the constructive renormalization-group problem, and
whose error is now explicit.

\hypertarget{the-decomposition}{%
\subsection{1. The decomposition}\label{the-decomposition}}

\href{mass-gap-obligations-lattice.md}{The obligations map} turns
``prove the mass gap'' into six statements about the Kogut--Susskind
Hamiltonian
\[H=\frac{\hbar c}{a}\Big[\frac{g^2}{2}\sum_{\ell}(-\Delta_\ell)
+\frac2{g^2}\sum_p\big(N-\operatorname{Re}\operatorname{tr}U_p\big)\Big],
\qquad \Delta_{a,L}=\frac{\hbar c}{a}\,\delta(g;N_s,G).\]

\begin{longtable}[]{@{}
  >{\raggedright\arraybackslash}p{(\columnwidth - 4\tabcolsep) * \real{0.3333}}
  >{\raggedright\arraybackslash}p{(\columnwidth - 4\tabcolsep) * \real{0.3333}}
  >{\raggedright\arraybackslash}p{(\columnwidth - 4\tabcolsep) * \real{0.3333}}@{}}
\toprule\noalign{}
\begin{minipage}[b]{\linewidth}\raggedright
\end{minipage} & \begin{minipage}[b]{\linewidth}\raggedright
statement
\end{minipage} & \begin{minipage}[b]{\linewidth}\raggedright
status
\end{minipage} \\
\midrule\noalign{}
\endhead
\bottomrule\noalign{}
\endlastfoot
T1 & finite lattice: unique physical ground state, \(\delta>0\) &
\textbf{proved here} \\
T2 & \(\delta\ge\gamma(g^2/2)C_2\) for \(g\ge g_0\), uniform in volume &
\textbf{proved here} from an imported theorem \\
T2\('\) & \(\liminf_{N_s}\delta>0\) for every \(g>0\) & open; false for
\(U(1)\) \\
T3 &
\(\delta_\infty(g)/(a\Lambda_{\rm lat}(g))\to m/(\hbar c\Lambda)\in(0,\infty)\)
& open \\
T4 & continuum infinite-volume theory with the axioms & finite-volume
ultraviolet stability known \\
S & small volume: \(\Delta=\delta_1g^{2/3}\hbar c/L\,[1+O(g^{2/3})]\) &
upper side proved, lower side conditional \\
\end{longtable}

The conjecture is T2\('\) together with T3, given T4.

\hypertarget{what-is-proved-here}{%
\subsection{2. What is proved here}\label{what-is-proved-here}}

\textbf{T1} (\href{mass-gap-obligations-lattice.md}{obligations map}
§2). For every \(a\), \(N_s\ge2\), \(g>0\) and compact connected \(G\):
\(H\) has compact resolvent, a simple strictly positive ground state, a
positive gap, and the ground state lies in the physical sector, which
inherits the gap. Ingredients: Peter--Weyl for \(-\Delta_\ell\), a
bounded perturbation argument for compact resolvent, Feynman--Kac
positivity, and invariance of the positive ground state under gauge
transformations.

\textbf{T2 in Hamiltonian form}
(\href{strong-coupling-uniform-gap.md}{T2 note}). The electric term is a
classical Hamiltonian in the Peter--Weyl partition with product ground
state and unit local gap after normalization, and the magnetic term is a
bounded translation-invariant perturbation of range \(\{0,1\}^3\) with
\[\beta=\frac{48N^2}{g^4(N^2-1)}\quad(SU(N)),\] so Yarotsky's
gap-stability theorem gives, for
\(g\ge g_0=(48N^2/[(N^2-1)\beta_*])^{1/4}\) and every lattice size,
\[\Delta_{a,L}\ \ge\ \gamma\,\frac{g^2}{2}\,\frac{N^2-1}{2N}\,\frac{\hbar c}{a},\]
with an infinite-volume ground state and exponential clustering.

\textbf{A Lieb--Robinson bound}
(\href{lieb-robinson-kogut-susskind.md}{note}). The unbounded part of
\(H\) is a sum of commuting single-link Laplacians, so the interaction
picture generated by it leaves a time-dependent bounded finite-range
interaction of unchanged norm and support, and the standard theorem
applies:
\[v\ \le\ \frac{32\,e\,Nc}{g^2},\qquad\text{uniform in the volume},\]
giving a strong-coupling correlation length
\(\xi\le C'aN/(\gamma C_2g^4)\), and showing that the lattice light cone
opens faster than \(c\) whenever \(g^2<32eN\).

\textbf{Upper bounds.} The Feynman--Bijl inequality for any
gauge-invariant \(f(U)\) (\href{lattice-gap-upper-bounds.md}{note}
Proposition 1); the strong-coupling gap pinned to order \(g^2\hbar c/a\)
on both sides (Corollary 3); the abelian gap bounded by the
plaquette-sine structure factor (Theorem 5), which is how \(U(1)\)
becomes gapless; the proof that \textbf{no operator linear in the
electric field with c-number coefficients is gauge invariant for a
non-abelian group} (Proposition 6), so the abelian gap-closing channel
is blocked by a string cost; and the averaged Polyakov-loop bound with
its necessary condition on the transverse susceptibility
(\href{polyakov-average-gap-bound.md}{note}).

\textbf{The finiteness clause reduced.} In the continuum,
\[m\ \le\ \frac{M_{k+1}}{M_k}\quad\text{for every }k,\qquad
M_k=\int E^k\,d\rho(E),\] with \(\rho\) the spectral measure of one
flowed zero-momentum correlator, and the ratios decrease as \(k\)
decreases (\href{moment-hierarchy-upper-bounds.md}{moment hierarchy}).
The best usable member,
\[m\ \le\ \frac{\hbar\,C_0(0)}{\int_0^\infty C_0(\tau)\,d\tau},\] needs
only T4 and \(C_0\not\equiv0\). The free-field value of the weakest
member is computed in closed form and checked twice
(\href{flowed-bound-free-field.md}{free-field note}):
\(m\le\frac{8}{3\sqrt\pi}\hbar c/\sqrt t\approx4.26\,\hbar c/\sqrt{8t}\),
with the coupling cancelling.

\textbf{Structure of the lower side.} With
\(\mathcal G=\{g:\liminf_{N_s}\delta>0\}\), T2\('\) is
\(\mathcal G=(0,\infty)\), which holds if and only if \(\mathcal G\) is
open and closed; T2 gives nonemptiness, the Lieb--Robinson bound removes
one obstacle to openness, and closedness is equivalent to the absence of
a zero-temperature bulk phase transition, hence, on its second-order
branch, to the uniqueness of the continuum limit
(\href{gapped-set-critical-coupling.md}{note}).

\hypertarget{what-has-been-ruled-out}{%
\subsection{3. What has been ruled out}\label{what-has-been-ruled-out}}

Four statements, each by computation.

\textbf{Perturbation around the free theory cannot give \(m<\infty\)
usefully.} The flowed bound is \(\approx4.26\,\hbar c/\sqrt{8t}\) in the
perturbative regime, weaker than the expected gap by
\(1/(\sqrt{8t}\,\Lambda)\gg1\); it reaches \(m\) only at
\(\sqrt{8t_*}\simeq4.26\,\hbar c/m\), where its coefficient stops being
perturbative (\href{flowed-bound-free-field.md}{free-field note} §7).

\textbf{Real-space blocking gains nothing.} Yarotsky's proof is a
polymer expansion, so the connected estimate a blocking step needs is
available; applying his criterion to a lattice of blocks of \(M^3\)
sites gives \[\beta_{\rm block}=\frac{24M^2N}{g^2\,\delta(g;M)},\qquad
\frac{\beta_{\rm block}}{\beta_{\rm direct}}\simeq\frac{3M^2}8,\]
because the block boundary grows like \(M^2\) while the block gap stays
flat at strong coupling and falls like \(1/M\) at small volume. Blocking
is strictly worse than no blocking at every coupling
(\href{blocking-criterion-monotone.md}{note}).

\textbf{Conjugation by the flow cannot help.} The gradient flow is a
diffeomorphism of the compact configuration space, so it is implemented
by a unitary and conjugation preserves the spectrum; the magnetic term
becomes the flowed action and the electric term becomes nonlocal with
range \(\sqrt{8t}\), and in the free case the two effects cancel mode by
mode (\(e^{-2tk^2}\) against \(e^{+2tk^2}\)). A locality-based criterion
appears to improve by \(e^{-4tk^2}\) only because it drops the locality
hypothesis, so the gain belongs to the truncation that restores it
(\href{flow-conjugation-truncation.md}{note}).

\textbf{No upper bound can give \(m>0\).} Finiteness of all negative
moments is necessary for a gap and insufficient:
\(d\rho=e^{-\hbar c\kappa/E}dE\) has every negative moment finite with
\(\inf\operatorname{supp}\rho=0\)
(\href{moment-hierarchy-upper-bounds.md}{moment hierarchy} §5).

\textbf{The small-volume corner is an ultraviolet problem.} The Schur
error of the reduction to the constant modes has a Berry part of
relative order \(g^2\log(\Lambda L)\), absorbed by the running coupling,
and a cubic part whose energy slope is \(g^2(\Lambda L)^2\), so the
bare-fiber reduction holds only for \(g\ll1/N_s\): a fixed-lattice
theorem, with the renormalized regime requiring a dressed vacuum
(\href{schur-error-ultraviolet.md}{Schur note}).

\hypertarget{where-the-non-abelian-structure-has-been-isolated}{%
\subsection{4. Where the non-abelian structure has been
isolated}\label{where-the-non-abelian-structure-has-been-isolated}}

Three statements distinguish the groups, and any proof of T2\('\) must
use one of them.

\begin{enumerate}
\def\labelenumi{\arabic{enumi}.}
\tightlist
\item
  \emph{The commutator potential.} In the zero-momentum sector the
  potential is \(\frac{g^2}2\sum_{i<j}|\vec x_i\times\vec x_j|^2\),
  identically zero for an abelian group; where it is present the
  transverse zero-point energy \(\hbar g|\vec x_i|/\sqrt m\) confines
  the valley and produces the gap \(\delta_1\hbar^{4/3}g^{2/3}m^{-2/3}\)
  (\href{low-dimensional-mass-gap.md}{G07} §3, claim C133).
\item
  \emph{Gauge invariance blocks the photon channel.} No operator linear
  in \(E\) with c-number coefficients commutes with Gauss's law for a
  non-abelian group, so the excitation through which the abelian gap
  closes must carry a Wilson line whose electric cost does not vanish at
  long wavelength (\href{lattice-gap-upper-bounds.md}{upper-bound note}
  Proposition 6).
\item
  \emph{The one-loop valley potential.} On the torus,
  \(U(a)=\frac{2\hbar c}{\pi^2L}\Phi(La)\) with
  \(\Phi(b)=\sum'_m(1-\cos m\cdot b)/|m|^4\), whose \(k=0\) term is
  exactly the four transverse oscillators of the matrix model; for
  compact \(U(1)\) the potential vanishes identically and the holonomy
  is free (\href{torus-valley-potential.md}{valley note} §3b).
\end{enumerate}

All three say that the abelian theory has flat directions that nothing
lifts.

\hypertarget{what-remains-and-of-what-kind}{%
\subsection{5. What remains, and of what
kind}\label{what-remains-and-of-what-kind}}

\begin{longtable}[]{@{}
  >{\raggedright\arraybackslash}p{(\columnwidth - 4\tabcolsep) * \real{0.3333}}
  >{\raggedright\arraybackslash}p{(\columnwidth - 4\tabcolsep) * \real{0.3333}}
  >{\raggedright\arraybackslash}p{(\columnwidth - 4\tabcolsep) * \real{0.3333}}@{}}
\toprule\noalign{}
\begin{minipage}[b]{\linewidth}\raggedright
obligation
\end{minipage} & \begin{minipage}[b]{\linewidth}\raggedright
kind
\end{minipage} & \begin{minipage}[b]{\linewidth}\raggedright
nearest tool
\end{minipage} \\
\midrule\noalign{}
\endhead
\bottomrule\noalign{}
\endlastfoot
T2\('\): closedness of \(\mathcal G\) & the conjecture & a change of
variables at each scale \\
openness of \(\mathcal G\) & technical & gap stability without
frustration-freeness \\
T3 & the conjecture, given T2\('\) & asymptotic-freedom control of
\(\delta_\infty(g)\) \\
T4 & construction & Balaban-type renormalization group \\
S lower side, renormalized & construction & dressed nonzero-mode
vacuum \\
(F1)-free finiteness & reduced & nontriviality of one flowed
correlator \\
\end{longtable}

The pattern across the notes is uniform: \textbf{every statement that is
uniform in the cutoff is a renormalization statement.} The
strong-coupling theorem is uniform in the volume at fixed cutoff; the
small-volume theorem is fixed-lattice; the flowed bounds are finite at
fixed lattice and need the flow's renormalization in the continuum. The
programme has pushed each of these to the point where the
renormalization step is the only thing left, and has shown that two ways
of avoiding it, blocking and perturbation around the free theory, do not
work.

\hypertarget{assessment}{%
\subsection{6. Assessment}\label{assessment}}

The route that survives is the one sketched at the outset and now made
precise: close the small-volume corner at fixed lattice; recognize the
blocked theory at each scale as a theory of the same form with a larger
coupling; iterate until the coupling reaches the strong-coupling
threshold, where T2 applies with a gap uniform in volume. The number of
steps is finite and logarithmic in the scale ratio, of order thirty for
\(SU(2)\) at a bare coupling \(g^2_{\rm UV}=1/2\)
(\href{blocking-step-obstruction.md}{blocking note} §5). The content of
each step is the change of variables, and that is exactly what has no
substitute: the notes show that projection alone loses, that the
connected estimate does not rescue it, and that the free theory is too
far away to expand around.

What this programme has added is a map with constants: six named
statements, two proved, the finiteness clause reduced to a single
correlator, two routes closed with explicit numbers, and three places
where the non-abelian structure is isolated in a form that a proof could
use. The remaining step is the one that has been open since the problem
was posed. \href{typical-field-strength-window.md}{The typical-field
note} locates it precisely: on typical configurations the flow
truncation is controlled at every coupling by lengthening the range as
\(\sqrt g\), and the whole obstruction is the supremum of the field
strength over the volume, whose control is probabilistic and therefore
Euclidean. The Hamiltonian route of this programme reaches its natural
limit there. \href{flow-jacobian-truncation-error.md}{The Jacobian note}
states the error: Gaussian in the truncation range over the flow radius,
times \(e^{2t\|G\|_\infty}\), so the step is controlled exactly where
the flowed field strength satisfies \(\|G\|_\infty\lesssim\ell^{-2}\).
The missing ingredient is a pointwise bound on that field strength,
where the flow supplies only an \(L^2\) bound.
\end{document}
